\documentclass[conference]{IEEEtran}
\IEEEoverridecommandlockouts
\usepackage{cite}
\usepackage{amsmath,amssymb,amsfonts}
\usepackage{graphicx}
\usepackage{textcomp}
\usepackage{xcolor}
\usepackage{booktabs}
\usepackage{url}
\usepackage{hyperref}

\newcommand{\todoplaceholder}[1]{\textcolor{red}{[PLACEHOLDER -- #1 -- fill in only after real experiments complete; do not fabricate]}}

\begin{document}

\title{ThreatGraph: Evaluating Knowledge Graph-Augmented RAG for Multi-Hop Cyber Threat Intelligence Question Answering}

\author{
\IEEEauthorblockN{Amber Hasan}
\IEEEauthorblockA{CS 6338 Knowledge Engineering}
}

\maketitle

\begin{abstract}
Cyber threat intelligence (CTI) analysts routinely need to answer questions
that require chaining several facts together -- for example, identifying
which threat groups use a given ATT\&CK technique, which malware those
groups employ to carry it out, and which mitigations defend against it.
Standard retrieval-augmented generation (RAG) systems retrieve semantically
similar text passages but have no explicit mechanism for multi-hop
relational reasoning. This paper presents ThreatGraph, a system that
augments text-only RAG with a MITRE ATT\&CK knowledge graph stored in Neo4j,
combining semantic chunk retrieval with bounded multi-hop graph traversal.
We built a 30-question multi-hop benchmark and a 60-triple extraction
benchmark, both derived programmatically from the official MITRE Enterprise
ATT\&CK v19.2 STIX data rather than hand-written, and designed a paired
evaluation comparing ThreatGraph against a text-only RAG baseline on the
same corpus, embedding model, generation model, and benchmark questions.
Across all 30 questions, ThreatGraph improved RAGAS Answer Relevancy from
0.000 to 0.430 and Context Recall from 0.067 to 0.644 (both paired
$p<10^{-5}$), with no significant change in Faithfulness and a
directionally positive but non-significant gain in Context Precision, and
no significant latency penalty despite retrieving 66.6\% more context
characters on average ($p=5.1\times10^{-30}$).
\end{abstract}

\begin{IEEEkeywords}
knowledge graphs, retrieval-augmented generation, cyber threat intelligence, MITRE ATT\&CK, question answering, GraphRAG
\end{IEEEkeywords}

\section{Introduction}
Cyber threat intelligence question answering frequently requires multi-hop
reasoning over structured relationships between threat actors, malware,
techniques, and mitigations -- relationships that MITRE ATT\&CK already
encodes explicitly as a graph. A question such as ``Which threat groups use
Process Injection (T1055), what malware associated with those groups uses
this technique, and which mitigations can help defend against it?'' requires
traversing \emph{ThreatGroup} $\rightarrow$ \emph{Software}
$\rightarrow$ \emph{Technique} $\leftarrow$ \emph{Mitigation} relationships,
not merely retrieving the single most semantically similar passage.
Standard RAG systems retrieve independently scored text chunks and have no
native mechanism for assembling such a chain. This motivates the central
research question: \emph{does adding MITRE ATT\&CK knowledge-graph
retrieval to RAG improve multi-hop CTI question answering compared with a
standard text-only RAG system, and if so, at what cost in latency and
irrelevant context?}

We do not assume the answer. GraphRAG approaches add retrieval complexity,
additional context, and additional latency; whether the added structure
translates into measurably better answers, worse answers (e.g. through
irrelevant graph facts diluting faithfulness), or no detectable difference
at this benchmark size is an empirical question this paper is designed to
measure, not assume.

\section{Related Work}
Retrieval-augmented generation (RAG) \cite{lewis2020} combines a
non-parametric retriever over an external corpus with a parametric
generator, grounding generated text in retrieved evidence rather than
relying solely on parameters learned at training time. Standard RAG
retrieves independently-scored passages and has no built-in mechanism for
combining facts that are individually retrievable but only answerable
together through an explicit relational chain -- precisely the gap
multi-hop CTI questions expose.

Graph-augmented RAG approaches address this by retrieving structured
relationships alongside or instead of free text. Microsoft Research's
GraphRAG \cite{edge2024graphrag} builds an LLM-derived entity-relationship
graph over a corpus, detects communities within it, and generates
hierarchical community summaries used at query time for
query-focused summarization -- a different mechanism (corpus-derived graph
construction plus summarization) from this project's approach (retrieval
over an existing authoritative knowledge graph, MITRE ATT\&CK, combined
with bounded multi-hop traversal), but motivated by the same observation
that pure vector retrieval under-serves questions requiring relational
reasoning across multiple entities.

MITRE ATT\&CK itself is maintained as a structured knowledge base (STIX
2.1) of adversary tactics, techniques, software, groups, and mitigations,
and is widely used as a reference ontology for cyber threat intelligence
tooling; this project uses it directly as the knowledge graph rather than
constructing one from unstructured text, reserving LLM-based extraction
(Section~\ref{sec:dataset}) for ingesting CTI sources that are not already
structured.

RAGAS \cite{es2024ragas} is used here for reference-free
(Faithfulness, Answer Relevancy) and reference-based (Context Precision,
Context Recall) automated RAG evaluation, as described in
Section~\ref{sec:results}.

\todoplaceholder{This section covers RAG, GraphRAG, MITRE ATT\&CK's role as
a structured CTI knowledge base, and RAGAS -- four load-bearing citations
whose core facts (authors, venue, year, core claims) the author of this
paper is confident are accurate. A stronger related-work section would add
specific prior work on (a) LLM-based cyber threat intelligence entity/
relationship extraction and (b) knowledge-graph question answering more
broadly; those are deliberately left out rather than filled with citations
this assistant could not verify against the original source, per the
project's instruction never to fabricate references.}

\section{System Architecture}
ThreatGraph consists of five stages, shown in Figure~\ref{fig:architecture}
(to be added): ingestion, knowledge extraction, entity resolution, graph
population, and dual retrieval-augmented generation.

\subsection{Ontology}
Entities are typed as \texttt{ThreatGroup}, \texttt{Technique},
\texttt{Software}, \texttt{Campaign}, \texttt{Mitigation},
\texttt{Vulnerability}, and \texttt{Document}. Relationships are restricted
to \texttt{USES}, \texttt{MITIGATES}, \texttt{ATTRIBUTED\_TO},
\texttt{EXPLOITS}, and \texttt{MENTIONED\_IN}, with a fixed set of allowed
(subject-type, predicate, object-type) triples enforced by a Pydantic schema
(\texttt{schemas.py}) shared across the extraction prompt, the extraction
validator, and the Neo4j schema, so the four representations cannot drift
out of sync.

\subsection{Ingestion and Knowledge Extraction}
The system ingests the official MITRE Enterprise ATT\&CK STIX bundle
(pinned to v19.2) directly for graph construction, since STIX already
provides structured, high-confidence entities and relationships. A separate
LLM-based extraction pipeline (prompt in
\texttt{prompts/extraction\_prompt.txt}) is implemented for ingesting
unstructured CTI text (e.g. CISA advisories) where no structured source
exists; its output is validated for shape and ontology conformance before
any entity resolution or graph write occurs, and evidence text is preserved
for every extracted relationship.

\subsection{Entity Resolution}
Extracted entity mentions are resolved against the canonical ATT\&CK entity
set in order of decreasing precision: (1) exact ATT\&CK identifier, (2)
known alias, (3) normalized canonical name, (4) embedding cosine similarity
above a configurable threshold as a last resort.

\subsection{Knowledge Graph}
Entities, relationships, text chunks, and provenance are stored in Neo4j
with uniqueness constraints on canonical identifiers, lookup indexes on
entity names, and a vector index over chunk embeddings. All writes are
parameterized, batched \texttt{MERGE} operations, making graph population
idempotent and safe to re-run.

\subsection{Dual RAG Systems}
System A (baseline) retrieves the top-$k$ semantically similar text chunks
via the Neo4j vector index and generates an answer from that context alone.
System B (ThreatGraph) performs the same vector retrieval, identifies seed
entities from ATT\&CK IDs mentioned in the question and from the top
retrieved chunks, performs bounded multi-hop graph traversal from those
seeds (capped at 3 hops and 15 facts), and combines text and
graph evidence into one context block before generation. Both systems share
the same embedding model, generation model, and prompt template, differing
only in the presence of graph retrieval.

For a hub entity with many edges of one type (e.g. a common technique with
dozens of \texttt{USES} edges from \texttt{Software} but only a few direct
\texttt{ThreatGroup} or \texttt{Mitigation} edges), naively truncating to the
first $k$ matched facts can crowd out rarer, often more directly relevant
edge types. ThreatGraph's traversal instead pools candidate facts per (hop
distance, entity type) bucket and round-robins across buckets, exhausting
1-hop facts before falling back to 2-hop facts, so one high-volume
relationship type cannot silently starve the context of ThreatGroup- or
Mitigation-level evidence -- a real failure mode found and fixed during
development (Section~\ref{sec:limitations}).

\section{Dataset and Benchmark Construction}
\label{sec:dataset}
Both benchmarks are derived programmatically from the official MITRE
Enterprise ATT\&CK v19.2 STIX bundle (\texttt{scripts/build\_benchmarks.py}),
not hand-authored, so every record traces back to a concrete STIX
relationship ID.

\texttt{gold\_triples.json} contains 60 \texttt{(subject, predicate,
object)} triples spanning \texttt{USES} and \texttt{MITIGATES}
relationships, used for extraction evaluation.

\texttt{eval\_queries.json} contains 30 multi-hop question-answer records
across three reasoning-path categories, summarized in
Table~\ref{tab:benchmark-categories}.

\begin{table}[h]
\centering
\caption{Multi-hop benchmark categories}
\label{tab:benchmark-categories}
\begin{tabular}{@{}lll@{}}
\toprule
Category & Hops & Count \\
\midrule
ThreatGroup$\to$Technique$\leftarrow$Mitigation & 2 & 14 \\
ThreatGroup$\to$Software$\to$Technique & 2 & 8 \\
ThreatGroup$\to$Software$\to$Technique$\leftarrow$Mitigation & 3 & 8 \\
\bottomrule
\end{tabular}
\end{table}

To prevent benchmark leakage, \texttt{expected\_answer} is withheld from
both RAG systems at answer time and attached to result records only after
generation, for scoring purposes.

\section{Experimental Setup}
Both systems share the following configuration, recorded automatically to
\texttt{results/experiment\_config.json} on every run for reproducibility:

\begin{itemize}
\item \textbf{Corpus}: MITRE Enterprise ATT\&CK v19.2 STIX, 1{,}798 active
  entities (176 ThreatGroup, 697 Technique, 825 Software, 56 Campaign, 44
  Mitigation) and 18{,}613 relationships, chunked into 3{,}079 text chunks.
\item \textbf{Embedding model}: \texttt{sentence-transformers/all-MiniLM-L6-v2},
  384 dimensions, run locally (no API key required for retrieval/indexing).
\item \textbf{Chunking}: 750-character chunks, 100-character overlap,
  character-based (not token-based) sliding window.
\item \textbf{Generation / extraction model}: \texttt{gemini-3.1-flash-lite}
  (temperature 0, max output tokens 2048, internal "thinking" disabled --
  see Section~\ref{sec:limitations}).
\item \textbf{RAGAS judge model}: \texttt{gpt-4o-mini} (temperature 0),
  separate from the generation/extraction model -- see
  Section~\ref{sec:limitations} for why.
\item \textbf{Retrieval budget}: top-5 text chunks for both systems; System B
  additionally retrieves up to 15 bounded graph facts within 3 hops.
\item \textbf{Benchmark}: 60 gold extraction triples, 30 multi-hop QA
  questions (14 two-hop ThreatGroup-Technique-Mitigation, 8 two-hop
  ThreatGroup-Software-Technique, 8 three-hop
  ThreatGroup-Software-Technique-Mitigation), all derived programmatically
  from the same STIX bundle (Section~\ref{sec:dataset}).
\end{itemize}

\section{Evaluation Methodology}
\subsection{Triple Extraction Metrics}
Precision, Recall, and F1 are computed by exact match on the canonical
triple $(\text{subject\_attack\_id}, \text{predicate}, \text{object\_attack\_id})$
after entity resolution. A true positive is a predicted triple present in
the gold set; a false positive is a predicted triple absent from the gold
set (including triples built from entities that could not be resolved to a
canonical ATT\&CK identifier); a false negative is a gold triple with no
matching predicted triple.
$$\text{Precision} = \frac{TP}{TP+FP}, \quad \text{Recall} = \frac{TP}{TP+FN}, \quad F1 = \frac{2 \cdot P \cdot R}{P+R}$$

\subsection{RAGAS Metrics}
Faithfulness and Answer Relevancy are \emph{reference-free}: Faithfulness
measures the fraction of claims in the generated answer that are entailed
by the retrieved context; Answer Relevancy measures how directly the answer
addresses the original question via back-generated question/original
question embedding similarity, independent of the retrieved context or the
gold answer. Context Precision and Context Recall are \emph{reference-based}:
Context Precision measures whether relevant retrieved evidence is ranked
near the top; Context Recall measures whether the retrieved context
contains the information needed to reconstruct the reference answer. We use
\texttt{ragas==0.2.15} (version pinning rationale in the project README) with
\texttt{gpt-4o-mini} as the judge LLM and a local
\texttt{sentence-transformers/all-MiniLM-L6-v2} embedding model for the
Answer Relevancy metric specifically (no embedding API key required).

\subsection{Human Evaluation}
A 1--5 rubric (Factuality, Completeness, Coherence) is defined for manual
rating of generated answers (\texttt{threatgraph/eval/human\_eval.py}).
\todoplaceholder{Human ratings, once actually collected -- do not report
scores here unless a human rater produced them.}

\subsection{Performance Instrumentation}
End-to-end, retrieval, and generation latency; input/output token counts;
estimated cost; and the number of retrieved chunks/graph facts are recorded
per question for both systems.

\subsection{Statistical Comparison}
Because both systems answer identical questions, all comparisons are
\emph{paired}: per-question differences, a paired $t$-test, and the
distribution-free Wilcoxon signed-rank test are reported together, with a
95\% confidence interval on the mean paired difference. With $n \le 30$
questions, a non-significant result is explicitly not interpreted as
evidence of equivalence.

\section{Results}
\label{sec:results}
All numbers below are real measurements from the 30-question benchmark and
the 60-triple extraction benchmark (results/comparison.csv,
results/ragas\_metrics.json, results/extraction\_metrics.json,
results/statistical\_analysis.json), using \texttt{gemini-3.1-flash-lite}
for extraction and answer generation (Section~\ref{sec:limitations}
discusses why this model, not the originally-intended
\texttt{gemini-3.8-flash}, was used) and \texttt{gpt-4o-mini} as the
separate RAGAS judge model (switched from Gemini after RAGAS's per-record
fan-out into several judge sub-calls each exhausted Gemini's free-tier
daily quota before a full pass could complete; see
Section~\ref{sec:limitations}).

\subsection{Triple Extraction}
\begin{table}[h]
\centering
\caption{Triple extraction performance (gemini-3.1-flash-lite, 60/60 gold-triple subjects attempted)}
\label{tab:extraction}
\begin{tabular}{@{}lr@{}}
\toprule
Metric & Value \\
\midrule
Precision (strict, vs. 60-triple sample) & 0.115 \\
Precision (adjusted, vs. full STIX ground truth) & 0.692 \\
Recall (vs. 60-triple sample) & 0.050 \\
F1 (strict) & 0.070 \\
True positives / predicted / gold & 3 / 26 / 60 \\
False positives: hallucinated & 8 \\
False positives: real fact outside sample & 15 \\
\bottomrule
\end{tabular}
\end{table}

The large gap between strict and adjusted precision is explained in
Section~\ref{sec:limitations}: most "false positives" under the strict
60-triple-sample definition are real STIX relationships the sample simply
did not include, not fabrications.

\subsection{Answer Quality (RAGAS)}
\begin{table}[h]
\centering
\caption{RAGAS scores, 30-question benchmark, 60 total judged answers (gpt-4o-mini judge; paired tests, n=30)}
\label{tab:results}
\begin{tabular}{@{}llll@{}}
\toprule
Metric & System A & System B & Paired $p$ (t / Wilcoxon) \\
\midrule
Faithfulness & 0.958 & 0.887 & 0.162 / 0.183 \\
Answer Relevancy & 0.000 & 0.430 & $3.4\times10^{-6}$ / $3.1\times10^{-5}$ \\
Context Precision & 0.296 & 0.415 & 0.068 / 0.098 \\
Context Recall & 0.067 & 0.644 & $1.1\times10^{-8}$ / $2.1\times10^{-5}$ \\
\bottomrule
\end{tabular}
\end{table}

Answer Relevancy and Context Recall both improve very significantly for
System B. Answer Relevancy rising from exactly 0.000 to 0.430 reflects a
qualitative pattern visible in the raw transcripts
(\texttt{results/raw\_experiment\_results.json}): System A very frequently
produces a variant of "the provided context does not contain information
linking this group to this technique," which RAGAS's answer-relevancy
metric (built on embedding similarity between the original question and
synthetic questions back-generated from the answer) scores as unrelated to
the question asked, since a refusal-to-answer is not *about* the same
specific entities as the question. Context Recall rising from 0.067 to
0.644 (+866\% relative, paired $t(29)=7.88$, $p=1.1\times10^{-8}$) means
System A's five text chunks very rarely happen to state the specific
technique-mitigation (or technique-software) pairing a question asks about,
while System B's added graph facts directly supply it. Context Precision is
directionally higher for System B (+40.5\% relative) but does not clear the
conventional $p<0.05$ threshold at $n=30$ (paired-$t$ $p=0.068$). Faithfulness
is not significantly different between systems ($p=0.162$); System A's
nominal edge is consistent with short refusal answers being trivially
faithful to whatever little context they do cite, rather than System B's
answers being less well-grounded in absolute terms.

\subsection{Latency and Context Volume}
\begin{table}[h]
\centering
\caption{Latency and context volume, 30-question benchmark (mean $\pm$ SD; paired tests, n=30)}
\label{tab:latency}
\begin{tabular}{@{}lll@{}}
\toprule
Metric & System A (Baseline) & System B (ThreatGraph) \\
\midrule
Mean total latency (s) & 5.49 $\pm$ 13.42 & 3.04 $\pm$ 10.16 \\
Mean context length (chars) & 1748.9 $\pm$ 564.4 & 2913.8 $\pm$ 598.9 \\
Retrieved text chunks & 5.0 (fixed) & 5.0 (fixed) \\
\bottomrule
\end{tabular}
\end{table}

Context length differs very significantly (paired $t(29)=51.3$, $p=5.1\times10^{-30}$;
Wilcoxon $p=1.7\times10^{-6}$): System B's graph facts add 1{,}164.8 characters
of context on average (+66.6\%), exactly as expected by construction -- this
is the intended experimental manipulation, not a finding in itself. Latency
is nominally lower for System B in this run (-44.6\% relative, mean
difference $-2.45$s) but the paired test is far from significant (paired-$t$
$p=0.238$; Wilcoxon $p=0.416$), and both systems show very high variance
(SD exceeding the mean), dominated by a small number of multi-second outlier
calls consistent with provider-side rate-limit backoff rather than a
systematic latency difference between the two retrieval strategies; see
Section~\ref{sec:limitations}. By reasoning depth, the 8 three-hop questions
show a larger (but still non-significant, $n=8$) latency reduction for
System B ($-75.3\%$) than the 22 two-hop questions ($-31.7\%$), the opposite
direction of what one might expect if graph traversal were adding
meaningful retrieval overhead.

\section{Discussion}

\textbf{Did graph retrieval improve multi-hop QA?} Yes, on the two metrics
that measure whether the system actually answers the question with the
right information: Answer Relevancy (0.000 $\to$ 0.430,
$p=3.4\times10^{-6}$) and Context Recall (0.067 $\to$ 0.644,
$p=1.1\times10^{-8}$). This is consistent with the research hypothesis --
MITRE ATT\&CK relationships needed to answer these questions (e.g. which
mitigations address a given technique) are structurally present in the
graph but not reliably recoverable from the top-5 semantically similar text
chunks alone, since a technique's mitigations are rarely co-located in the
same passage as the technique's description.

\textbf{Did the improvement differ between 2-hop and 3-hop questions?} Yes,
and in an unexpected direction. For the 22 two-hop questions, System B's
mean Answer Relevancy is 0.586 vs. System A's 0.000. For the 8 three-hop
questions, System B's mean Answer Relevancy is also 0.000 -- identical to
the baseline -- even though Context Recall still improves for three-hop
questions (0.000 $\to$ 0.275). In other words, for three-hop questions the
graph retrieval layer successfully supplies more of the needed facts into
the context, but the generation step does not reliably turn that richer
context into an answer RAGAS judges as relevant to the specific
multi-entity question asked. This points to a synthesis bottleneck rather
than a retrieval bottleneck for the hardest reasoning-depth category, and
is plausibly specific to \texttt{gemini-3.1-flash-lite}'s capability at
synthesizing several chained facts (group $\to$ software $\to$ technique
$\to$ mitigation) into a coherent answer, not necessarily a property of
graph-augmented retrieval in general.

\textbf{Did graph retrieval introduce irrelevant context?} Context
Precision is directionally higher for System B (0.296 $\to$ 0.415,
+40.5\%) but does not reach conventional significance at $n=30$
($p=0.068$), so the data here does not support a strong claim either way.
Qualitatively, the round-robin bucketing fix described in
Section~\ref{sec:limitations} was built specifically to prevent one
high-volume relationship type from drowning out rarer, more relevant graph
facts; its effect on Context Precision specifically was not isolated as a
separate ablation in this evaluation and would be a natural next
experiment.

\textbf{Did graph retrieval increase latency?} No measurable systematic
increase: retrieval itself (vector search plus bounded graph traversal) is
under 0.1s on average for both systems, and total latency -- dominated by
LLM generation, which both systems share -- showed no significant paired
difference (Table~\ref{tab:latency}, $p=0.238$). The additional ~1,165
characters of graph-fact context did not translate into materially slower
generation in this evaluation.

\section{Limitations}
\label{sec:limitations}
This evaluation uses a 30-question benchmark, which limits statistical
power for the paired significance tests reported in
Section~\ref{sec:results}; non-significant results should
not be read as equivalence. Chunking is character-based rather than
token-based. The local embedding model
(\texttt{sentence-transformers/all-MiniLM-L6-v2}, 384 dimensions) is smaller
than commonly used hosted embedding models, which may affect retrieval
quality for both systems equally. Graph retrieval currently covers two
reasoning-path shapes derived from \texttt{USES}/\texttt{MITIGATES}
relationships only; \texttt{EXPLOITS} and \texttt{ATTRIBUTED\_TO} paths are
supported by the ontology and bounded-traversal code but are not represented
in the current benchmark.

\textbf{Model/provider constraints.} The free tier of the LLM provider used
for this evaluation (Google Gemini) imposes per-model request quotas. An
initial run with \texttt{gemini-3.8-flash} was blocked by a 20-requests/day
quota after exploratory testing alone consumed most of it; that model also
spent most of its output-token budget on internal reasoning ("thinking")
tokens unless that behavior was explicitly disabled, which otherwise
truncated structured-JSON extraction output mid-string. The reported generation/extraction
experiments instead use \texttt{gemini-3.1-flash-lite}, which has a
separate, more usable per-minute quota (15 requests/minute) but is a
smaller, lower-capacity model; both the latency and answer-quality results
in this paper should be read as specific to this model, not as a general
claim about GraphRAG vs. baseline latency or quality under a larger
generation model. RAGAS judging hit a further, distinct quota wall even on
\texttt{gemini-3.1-flash-lite}: a single RAGAS metric does not correspond
to one LLM call per record -- Faithfulness alone issues one call to extract
claims plus one call per claim to verify it, and Context Precision/Recall
issue calls per retrieved context item -- so evaluating 60 records across 4
metrics fanned out into several hundred real API calls and exhausted
Gemini's 500-requests/day free-tier limit for that model before a full pass
completed, independent of and in addition to the per-minute limit above.
The RAGAS judge was switched to \texttt{gpt-4o-mini} (a separate provider
and quota, paid but inexpensive -- an estimated well under \$1 for this
evaluation's full 60-record, 4-metric pass, though this project did not
verify the exact billed amount against OpenAI's dashboard) to obtain
complete, reliable judged scores;
the practical lesson is that RAGAS's per-metric call-fan-out should be
budgeted for explicitly against a provider's actual daily/per-minute
limits, not approximated as one call per (record, metric) pair.

\textbf{Graph-traversal correctness bugs found during development.}
Development surfaced and fixed two bugs in the bounded graph-traversal
fallback used when the Neo4j instance lacks the APOC plugin: an undirected
neighbor query that reported relationship direction based on which side
matched the seed node rather than the edge's actual stored direction
(silently inverting facts like ``Software USES Technique'' into ``Technique
USES Software''), and a result-key mismatch that caused almost all retrieved
facts to collapse together during deduplication. Both were caught by
manually inspecting retrieved graph facts against known ATT\&CK relationships
during development, not by an automated check.

\textbf{Extraction evaluation findings.} Running real LLM extraction against
60 gold triples (Table~\ref{tab:extraction}) surfaced three specific,
addressable issues rather than a single aggregate quality score. First, MITRE
ATT\&CK's STIX descriptions embed inline bibliographic citation markers (for
example ``(Citation: Kaspersky Winnti April 2013)''); the initial extraction
prompt had no instruction to ignore these, so the model frequently extracted
them as spurious \texttt{Document} entities, driving a high rate of
unresolvable entities. Adding an explicit rule to ignore citation markers,
and an explicit table of every allowed (subject-type, predicate,
object-type) shape rather than a bare predicate list, reduced ontology-shape
rejections from 46 to 2 out of 71 relationships attempted on the same
20-entity re-test. Second, the model sometimes inverted the intended
direction of \texttt{ATTRIBUTED\_TO} (extracting \texttt{ThreatGroup
ATTRIBUTED\_TO Campaign} where this project's schema defines \texttt{Campaign
ATTRIBUTED\_TO ThreatGroup}) -- both directions are a defensible modeling
choice, but the mismatch between the model's assumption and this schema's
convention is itself a finding worth noting for anyone reusing this
ontology. Third, and most significant for interpreting Table~\ref{tab:extraction}:
because \texttt{gold\_triples.json} is a 60-triple \emph{sample} of MITRE
ATT\&CK's full relationship set rather than an exhaustive list, a meaningful
fraction of "false positives" under the strict sample-only definition are
real, verifiable STIX relationships that simply were not drawn into the
sample -- in the full run, 15 of 23 false positives (65\%) were verified
against the complete parsed STIX relationship set and are factually correct;
only 8 were genuine fabrications. Crediting the 15 sample-external-but-real
triples raises precision from 0.115 to an adjusted 0.692
(\texttt{threatgraph.eval.extraction\_metrics.classify\_false\_positives}).
Recall (0.05) is not similarly adjustable, since \texttt{gold\_triples.json}
is the intended recall target by construction; the low recall reflects that
a single description text often does not happen to state the one specific
relationship a diversity-balanced sampling procedure selected as that
entity's gold triple, which is a property of this evaluation design, not
necessarily of extraction quality.

\section{Conclusion}
On this 30-question benchmark, adding MITRE ATT\&CK knowledge-graph
retrieval to RAG produced a statistically significant improvement in
whether the generated answer actually addresses the specific multi-hop
question asked (Answer Relevancy, $p=3.4\times10^{-6}$) and in whether the
retrieved evidence contains what is needed to support the reference answer
(Context Recall, $p=1.1\times10^{-8}$), with no significant cost to
Faithfulness or latency, and a non-significant directional improvement in
Context Precision. This improvement was concentrated in two-hop questions;
for the harder three-hop category, ThreatGraph retrieved more of the
necessary evidence (Context Recall still improved) but this did not
translate into a measurable Answer Relevancy gain, suggesting the
bottleneck for the hardest questions in this evaluation was answer
synthesis rather than retrieval. These results support the research
hypothesis for the metrics and question types where it was tested, while
leaving open how much of the effect is attributable to the specific
(small, free-tier) generation model used rather than graph-augmented
retrieval as a general technique. The extraction and graph-construction
pipeline is independently useful and genuinely validated: real MITRE
ATT\&CK STIX data populates a working Neo4j knowledge graph of 1,798
entities and 18,613 relationships, and the development process itself
surfaced and fixed multiple concrete correctness bugs (Section~\ref{sec:limitations})
in ways that should generalize to other GraphRAG implementations over
similar bounded-traversal patterns.

\section*{Acknowledgment}
This project makes use of MITRE ATT\&CK\textregistered{}.

\begin{thebibliography}{00}
\bibitem{lewis2020} P. Lewis et al., ``Retrieval-Augmented Generation for Knowledge-Intensive NLP Tasks,'' \emph{Advances in Neural Information Processing Systems}, vol. 33, 2020.
\bibitem{es2024ragas} S. Es, J. James, L. Espinosa-Anke, and S. Schockaert, ``RAGAS: Automated Evaluation of Retrieval Augmented Generation,'' \emph{Proceedings of the 18th Conference of the European Chapter of the Association for Computational Linguistics: System Demonstrations}, 2024.
\bibitem{edge2024graphrag} D. Edge, H. Trinh, N. Cheng, J. Bradley, A. Chao, A. Mody, S. Truitt, and J. Larson, ``From Local to Global: A Graph RAG Approach to Query-Focused Summarization,'' \emph{arXiv preprint arXiv:2404.16130}, 2024.
\bibitem{mitreattack} MITRE, ``MITRE ATT\&CK,'' \url{https://attack.mitre.org/}, accessed 2026.
\todoplaceholder{Add remaining references only after verifying each citation against its real source; do not leave fabricated entries here.}
\end{thebibliography}

\end{document}
